
\documentclass{article}
\usepackage{colortbl}
\usepackage{makecell}
\usepackage{multirow}
\usepackage{supertabular}

\begin{document}

\newcounter{utterance}

\centering \large Interaction Transcript for game `guesswhat', experiment `Abs\_Level\_2', episode 6 with Qwen3{-}30B{-}A3B{-}t0.0.
\vspace{24pt}

{ \footnotesize  \setcounter{utterance}{1}
\setlength{\tabcolsep}{0pt}
\begin{supertabular}{c@{$\;$}|p{.15\linewidth}@{}p{.15\linewidth}p{.15\linewidth}p{.15\linewidth}p{.15\linewidth}p{.15\linewidth}}
    \# & $\;$A & \multicolumn{4}{c}{Game Master} & $\;\:$B\\
    \hline

    \theutterance \stepcounter{utterance}  
    & & \multicolumn{4}{p{0.6\linewidth}}{
        \cellcolor[rgb]{0.9,0.9,0.9}{
            \makecell[{{p{\linewidth}}}]{
                \texttt{\tiny{[P1$\langle$GM]}}
                \texttt{Let's play a guessing game. Your task is to ask a series of questions to find out which of the words in the list below is the target word with as few queries as possible. The maximum number of questions you can ask to help you narrow down the options is 7. Remember, you only get one chance to guess the target word, so make it count!} \\
\\ 
\texttt{The rules of the game are: } \\
\\ 
\texttt{1. Only ask one question per turn. DO NOT ASK MORE THAN ONE QUESTION.} \\
\\ 
\texttt{2. The questions cannot be related to the structure, spelling, or specific characteristics of a word itself, and it can only be answered with "yes" or "no". } \\
\\ 
\texttt{3. Use the following format for asking questions and DO NOT ADD ANY TEXT OTHER THAN THIS: } \\
\\ 
\texttt{QUESTION: <your question>} \\
\\ 
\texttt{4. You have only one opportunity to guess the target word. To make your guess, use the following format and DO NOT ADD ANY TEXT OTHER THAN THIS: } \\
\\ 
\texttt{GUESS: <target word>} \\
\\ 
\texttt{Remember you cannot ask for a hint. Make your question brief. Let’s start!} \\
\\ 
\texttt{List of candidate words: {[}'tsar', 'king', 'princes', 'prince', 'noon', 'dusk', 'afternoon', 'evening'{]}} \\
            }
        }
    }
    & \\ \\

    \theutterance \stepcounter{utterance}  
    & \multicolumn{4}{p{0.6\linewidth}}{
        \cellcolor[rgb]{0.9,0.9,0.9}{
            \makecell[{{p{\linewidth}}}]{
                \texttt{\tiny{[P1$\rangle$GM]}}
                \texttt{<think>} \\
\texttt{Okay, let's see. I need to figure out the target word from the list using as few yes/no questions as possible, max 7. The words are 'tsar', 'king', 'princes', 'prince', 'noon', 'dusk', 'afternoon', 'evening'. } \\
\\ 
\texttt{First, I should think about categories that can split the list into smaller groups. Since the questions can't relate to spelling or structure, I have to think about meanings or contexts.} \\
\\ 
\texttt{Maybe start with something like "Is the word related to a time of day?" That would separate 'noon', 'dusk', 'afternoon', 'evening' from the others. But wait, 'tsar', 'king', 'princes', 'prince' are titles. So if the answer is yes, then the target is in the time group. If no, then it's a title. But I need to check if that's a valid question. The question is about the word's meaning, not structure, so that's allowed.} \\
\\ 
\texttt{But maybe there's a better first question. Let's think. Another approach: "Is the word a single person?" Because 'princes' and 'prince' are related to people, but 'prince' is singular. Wait, but 'princes' is plural. 'King' and 'tsar' are also single people. So maybe that's not the best. Alternatively, "Is the word a} \\
            }
        }
    }
    & & \\ \\

    \theutterance \stepcounter{utterance}  
    & & & \multicolumn{2}{p{0.3\linewidth}}{
        \cellcolor[rgb]{0.9,0.9,0.9}{
            \makecell[{{p{\linewidth}}}]{
                \texttt{\tiny{[GM$|$GM]}}
                \texttt{Invalid format. Guesser must use the form 'QUESTION: ' or 'GUESS: '.} \\
            }
        }
    }
    & & \\ \\

    \theutterance \stepcounter{utterance}  
    & & & \multicolumn{2}{p{0.3\linewidth}}{
        \cellcolor[rgb]{0.9,0.9,0.9}{
            \makecell[{{p{\linewidth}}}]{
                \texttt{\tiny{[GM$|$GM]}}
                \texttt{abort game} \\
            }
        }
    }
    & & \\ \\

\end{supertabular}
}

\end{document}
